\thispagestyle{empty}

\begin{center}
\textbf{Министерство науки и высшего образования Российской Федерации}
\\
{\fontsize{9}{10.8}\selectfont
ФЕДЕРАЛЬНОЕ ГОСУДАРСТВЕННОЕ АВТОНОМНОЕ ОБРАЗОВАТЕЛЬНОЕ УЧРЕЖДЕНИЕ ВЫСШЕГО ОБРАЗОВАНИЯ}\\
\textbf{«Национальный исследовательский университет ИТМО»\\
(Университет ИТМО)}\\[1em]
\end{center}

\begin{flushleft}
\textbf{Факультет} \hspace{1em} \textbf{ПИиКТ}\\[0.5em]
\textbf{Образовательная программа} \textbf{Веб-технологии}\\[1em]

\textbf{Направление подготовки (специальность)} \textbf{09.04.04 Программная инженерия}\\[2em]
\end{flushleft}
\begin{center}
{\fontsize{16}{19.2}\selectfont ОТЧЕТ}\\[0.5em]
\mbox{по практике производственной, технологической (проектно-технологической)}\\[2em]
\end{center}

\begin{flushleft}

Тема задания: \textit{Реализация кросс-языковой браузерной NLP-системы анализа тональности}\\[1em]

Обучающийся: \textit{Юсеф Ая, P4207}\\[2em]

Согласовано:\\[0.5em]
Руководитель практики от университета: \textit{Государев И.Б., доцент, кандидат педагогических наук, факультет ПИиКТ}\\[2em]
\end{flushleft}

\begin{flushright}
\begin{tabular}{l l}
Практика пройдена с оценкой & \hspace{1em}\\
\\
Дата \textbf{29.04.2026} & \hspace{1em}
\end{tabular}
\end{flushright}

\vfill

\begin{center}
    Санкт-Петербург\\2026
\end{center}